\documentclass{ximera}
%herschreven 13/11/26
\title{Herhaling: eenheden en standaardvorm}
\author{Daan Lipkens}

\begin{document}
\begin{abstract}
    Het SI-stelsel, eenheden omzetten, beduidende cijfers en de standaardvorm.
\end{abstract}
\maketitle

\section*{Inleiding}\label{hoofdstuk:standaardvorm}
In 2023 berichtte een krant over een meteoriet die ``zo groot als een corgi'' was (figuur~\ref{fig:corgi}). Hoe groot is dat? Een puppy of een volwassen hond, een mannetje of een vrouwtje: het maakt een groot verschil. Een getal zonder duidelijke eenheid zegt weinig. Daarom spreken wetenschappers overal ter wereld dezelfde eenheden af: het \textit{Système international d'unités}, kortweg het SI-stelsel.

% AFBEELDING: corgi
% bron: https://commons.wikimedia.org/wiki/File:Welsh_Pembroke_Corgi.jpg
\begin{figure}[ht!]
    \captionsetup{width=0.85\textwidth}
    \begin{center}
        \includegraphics[width=0.45\textwidth]{pic5/appendix/corgi.jpg}
        \fotobron{Dog breed facts}{CC BY-SA 4.0}{https://commons.wikimedia.org/wiki/File:Welsh_Pembroke_Corgi.jpg}
    \end{center}
    \caption{Een Pembroke Welsh corgi. Maar hoe groot is ``zo groot als een corgi'' nu precies?}
    \label{fig:corgi}
\end{figure}

Ook de \emph{nauwkeurigheid} van een meetresultaat moet je kunnen aflezen. Een lengte van $2 \, \mathrm{m}$ is iets anders dan een lengte van $2{,}000 \, \mathrm{m}$. In deze appendix zetten we de afspraken op een rij die je in elke oefening van deze cursus gebruikt:
\begin{itemize}
    \item het SI-stelsel en de prefixen (sectie~\verwijs[Het SI-stelsel en prefixen]{sec:si_stelsel}{appendix_standaardvorm_si});
    \item eenheden omzetten: oppervlakte, volume, samengestelde eenheden en tijd (sectie~\verwijs[Omzetten van eenheden]{sec:omzetten_oppervlakte_volume}{appendix_standaardvorm_omzetten_eenheden});
    \item beduidende cijfers, wetenschappelijke notatie en standaardvorm (sectie~\verwijs[Beduidende cijfers en standaardvorm]{sec:beduidende_cijfers}{appendix_standaardvorm_bc}).
\end{itemize}

\begin{definition}[title={Wat je na deze appendix moet kennen en kunnen}]\nl
    \begin{itemize}
        \item de zeven SI-basiseenheden en de prefixen van giga tot nano gebruiken;
        \item lengte, oppervlakte, volume, samengestelde eenheden en tijd omzetten met machten van 10;
        \item het aantal beduidende cijfers van een meetresultaat bepalen en een resultaat correct afronden na een som, verschil, product of quotiënt;
        \item een eindantwoord in standaardvorm noteren.
    \end{itemize}
\end{definition}

\end{document}
